% !TeX program = lualatex
% !TeX encoding = utf8
% !TeX spellcheck = uk_UA
 !TEX indexprogram = upmendex
% !TeX root =./ElMagConspect.tex

%=========================================================
\chapter*{Передмова}\label{\currfilebase}
\addcontentsline{toc}{chapter}{Передмова}
%=========================================================


Пей конспект лекцій охоплює курс «Електрика та магнетизм», який
читається на 2-му курсі для здобувачів ступеня бакалавра спеціальності
E6~«Прикладна фізика та наноматеріали». Матеріал побудовано за
класичною схемою: від електростатики у вакуумі, провідниках і
діелектриках --- через постійний струм і магнітне поле --- до
рівнянь Максвелла, електромагнітних хвиль та коливальних
процесів у колах і хвилеводах.

Кожен розділ
супроводжується короткою історичною чи філософською довідкою
яка нагадує, що за рівняннями стоїть жива історія експериментальної думки. Там,
де це можливо, наведено фізичні аналогії (наприклад, зв'язок
еквіпотенціальних поверхонь з рельєфом місцевості), які
допомагають побудувати інтуїцію, перш ніж переходити до
формального апарату.

Текст має наскрізну систему позначок, яка полегшує
самостійну роботу з конспектом:
\begin{itemize}
    \item знаком <<!>> виділено принципові зауваження та
    поширені помилки, на які варто звернути особливу увагу;
    \item знаком <<?>> позначено задачі для самостійного
    розв'язання, вбудовані безпосередньо в текст лекції, одразу
    після відповідного теоретичного матеріалу;
    \item зірочкою <<*>> відмічено додатковий матеріал —
    складніші виведення чи приклади, які виходять за межі
    обов'язкового мінімуму, але поглиблюють розуміння теми;
    \item наприкінці кожного розділу наведено список контрольних
    питань для підготовки до семестрового контролю.
\end{itemize}

Виклад послідовно спирається на гаусову систему одиниць,
природну для теоретичної електродинаміки, з переходом до СІ
там, де це методично виправдано (додаток~А містить таблицю
переведення між системами).

Конспект адресований насамперед студентам фізико-технічного
інституту, які вивчають електрику та магнетизм у межах
дисципліни «Електрика та магнетизм», але може бути корисним
і всім, хто прагне поєднати строгий математичний виклад
класичної електродинаміки з живим фізичним розумінням явищ,
що стоять за формулами.